\clearpage
\begingroup
\fancyhead[L]{\small NACHLASS UND BRIEFWECHSEL}
\renewcommand{\headrulewidth}{0pt}
\par\medskip
\Needspace{8\baselineskip}
\pdfbookmark[2]{[3]}{pallas-IX-fragment3}
\begin{center}[3.]\end{center}
\begingroup
\small
\newcommand{\IXmars}{\text{\AstroSymbols ♂}}
\newcommand{\IXpallas}{\text{\AstroSymbols ⚴}}
\par\noindent\hspace*{\parindent}[Durch Multiplication mit den Factoren der Gleichungen 2) erhält man endlich die Ausdrücke für
\(\frac{\mathrm{d}\log\mathrm{hyp}\,n}{n\,\mathrm{d}t}\) und \(\frac{\mathrm{d}\varphi}{n\,\mathrm{d}t}\); der erstere ergibt, wenn die Integration eben so wie Seite 523---524 ausgeführt wird,
die Störungen des Logarithmen der halben grossen Axe und die des ersten Theils der Epoche, der zweite
die von \(\varphi\). Es fand sich:
\[
\begin{aligned}
\delta\log a&=1\sin(\IXmars-3\IXpallas+10^0\ 27'\ 49'')
&\delta\!\int n\,\mathrm{d}t&=0\overset{\prime\prime}{,}18\sin(\IXmars-3\IXpallas+280^0\ 27'\ 49'')\\
&\quad+1\sin(\IXmars-2\IXpallas+202\ 20\ 18)
&&\quad+0{,}16\sin(\IXmars-2\IXpallas+292\ 20\ 18)\\
&\quad+6\sin(\IXmars-\IXpallas+90\ 54\ 56)
&&\quad+0{,}28\sin(\IXmars-\IXpallas+180\ 54\ 56)\\
&\quad+11\sin(2\IXmars-5\IXpallas+124\ 17\ 38)
&&\quad+8{,}14\sin(2\IXmars-5\IXpallas+34\ 17\ 38)\\
&\quad+2\sin(2\IXmars-4\IXpallas+162\ 24\ 56)
&&\quad+0{,}12\sin(2\IXmars-4\IXpallas+252\ 24\ 56)\\
&\quad+2\sin(3\IXmars-7\IXpallas+76\ 26\ 28)
&&\quad+0{,}38\sin(3\IXmars-7\IXpallas+166\ 26\ 28)\\
&\quad+2\sin(4\IXmars-10\IXpallas+171\ 18\ 55)
&&\quad+0{,}79\sin(4\IXmars-10\IXpallas+81\ 18\ 55)\\
&\quad+1\sin(5\IXmars-12\IXpallas+120\ 44\ 45)
&&\quad+0{,}28\sin(5\IXmars-12\IXpallas+210\ 44\ 45).
\end{aligned}
\]
\[
\begin{aligned}
\delta\varphi&=0\overset{\prime\prime}{,}15\sin(\IXmars+335^0\ 22'\ 13'')\\
&\quad+0{,}49\sin(2\IXmars-5\IXpallas+106\ 33\ 27).
\end{aligned}
\]
\begin{center}
\begin{tabular}{@{}lrl@{}}
Säcularer Theil von \(\delta\varphi\): & täglich & \(-0\overset{\prime\prime}{,}000\,0059\)\\
& jährlich & \(-0\overset{\prime\prime}{,}00215\).\\
\end{tabular}
\end{center}
\par\noindent\hspace*{\parindent}\(\delta\log a\) ist in Einheiten der 7. Decimale gegeben, und es sind alle Gleichungen berücksichtigt, welche
den Betrag einer solchen Einheit resp.\ den von \(0\overset{\prime\prime}{,}10\) erreichen.]
\par\smallskip\noindent VII.\hfill 75\par
\par\endgroup

\par\medskip
\SourceLeaf{600}
\Needspace{8\baselineskip}
\pdfbookmark[2]{[4]}{pallas-IX-fragment4}
\begin{center}[4.]\end{center}
\begingroup
\small
\par\noindent\hspace*{\parindent}[Die Reihen, in welche im Vorigen die Grössen \(\frac{x'-x}{\rho^3}\), \(\frac{y'-y}{\rho^3}\), \(\frac{z'-z}{\rho^3}\) entwickelt worden sind,
haben eine so langsame Convergenz, dass trotz der Kleinheit der einzelnen Störungsgleichungen 24 Marsörter
und eben so viele Pallasörter der Rechnung zum Grunde gelegt werden mussten. Man kann aber hier
auch anders verfahren, indem man nemlich die zu entwickelnden Grössen mit einem Factor \(\lambda\) multiplicirt,
der so beschaffen ist, dass einerseits die Entwickelung der Producte \(\frac{x'-x}{\rho^3}\lambda\) etc. möglichst gut convergirt,
während andererseits die Entwickelung von \(\lambda^{-1}\) zwar langsam convergirt, sich aber leicht mit Hülfe analytischer
Formeln beliebig weit fortsetzen lässt. Man kann dann die Producte \(\frac{x'-x}{\rho^3}\lambda\) etc. interpolatorisch
leicht entwickeln, indem man nur wenige Marsörter zum Grunde legt, und die Entwickelung von \(\lambda^{-1}\) mit
wenig Mühe so weit fortsetzen, wie es nöthig scheint. Die Multiplication beider Entwickelungen gibt hierauf
die gesuchte Entwickelung von \(\frac{x'-x}{\rho^3}\) etc. Auch kann diese Methode zur Controlle der auf rein interpolatorischem
Wege gewonnenen Entwickelungen dienen.
\par\noindent\hspace*{\parindent}Der Factor \(\lambda\) wurde auf folgende Weise bestimmt:]
\par\endgroup
\clearpage
\endgroup
